\documentclass[10pt]{article}

\usepackage[utf8]{inputenc}

\usepackage[T1]{fontenc}

\usepackage{amsmath}

\usepackage{amsfonts}

\usepackage{amssymb}

\usepackage[version=4]{mhchem}

\usepackage{stmaryrd}

\usepackage{bbold}

\usepackage{mathrsfs}



\begin{document}

РАВНОРАСПРЕДЕЛЕНИЯ ЭНЕРГИИ ПРИНЦИП: дЛЯ классических систем с квадратичным по фазовым переменным гамильтонианом (идеальный газ) каждому квадратичному члену соответствует вклад в энергию (в расчете на одну частицу), равный $\theta / 2$, где $\theta$ - температура в энергетических единицах. В частности, для одноатомного идеального трехмерного газа число таких квадратичных членов $l=3$ и полная энергия $\varepsilon=3 \theta / 2$. Для нелинейной $n$-атомной молекулы $l=6(n-1)$.



Лlum.: [1] Л а н д а у Л. Д., Л и ф ш и ц Е. М., Статистическая физика, ч. 1, М., $1976 . \quad$ Д. П. Санкович. РАДИАЛЬНАЯ ФУНКЦИЯ РАСПРЕДЕЛЕНИЯ - см. Корреляционная функция.



РАДИАЛЬНОЕ УРАВНЕНИЕ ШРЕДИНГЕРА - ДВУХЧастичное Шредингера уравнение



\[

\left(-\Delta_{x}+v(x)-k^{2}\right) \Psi(x, k)=0,

\]



$\boldsymbol{x}, \boldsymbol{k} \in \mathbb{R}^{3}$, пля сферически симметричного потенщиала $v(\boldsymbol{x})=v(|\boldsymbol{x}|)$, приводящееся к системе независимых уравнений



\[

-\frac{d^{2} u_{l}}{d r^{2}}+\left[v(r)+\frac{l(l+1)}{r^{2}}-k^{2}\right] u_{l}=0, \quad l=0,1,2, \ldots,

\]



$\boldsymbol{r}=|\boldsymbol{x}|$. Функщии $u_{l}=u_{l}(\boldsymbol{r}, \boldsymbol{k})$ являются коэффициентами в разложении функции $\Psi(x, k)$ в ряд по многочленам Лежандра:



\[

\Psi(\boldsymbol{x}, \boldsymbol{k})=\sum_{l=0}^{\infty}(2 l+1) i^{l} \frac{u_{l}(\boldsymbol{r}, \boldsymbol{k})}{k r} P_{l}(\cos \theta)

\]



где $\cos \theta=(\boldsymbol{x}, \boldsymbol{k}) /(|\boldsymbol{k}||\boldsymbol{x}|)$. Решения задачи рассеяния в этом случае определяются как решения уравнения (1) с асимптотикой при $r \rightarrow \infty$ вида



\[

u_{l}(r, k) \sim \exp \left\{i \delta_{l}(k)\right\} \sin \left(k r-\frac{l \pi}{2}+\delta_{l}(k)\right) .

\]



Уравнения (1) называются радиальными уравне ниями Ш редингера, функции $u_{l}(\boldsymbol{r}, \boldsymbol{k})-$ парци альными компонентами волновой функции $\Psi(\boldsymbol{x}, \boldsymbol{k})$, а функции $\delta_{l}(\boldsymbol{k})$ - парциальны м и фазам и рассея ния.



Jum.: [1] Faxe n H., Holtzmark J., "Z. Phys.», 1927, Bd 45, S. 307-24; [2] Л а н да у Л. Д., Л и ф ш и ц Е. М., Квантовая механика. Нерелятивистская теория, 4 изд., М., 1989.



С.П. Меркурьев, С. Л. Яковлев.



PAДИАЧИОННАЯ ШИРИНА УРОВНЯ - см. Ширина уровня.



PАДИАЦИОННЫЕ ПОПРАВКИ в к в а н т о в о й э Л е кт р о д и н а м и к е - поправки к значениям некоторых физических величин и к сечениям различных процессов (вычисленным по формулам релятивистской квантовой механики), обусловленные взаимодействием заряженной частицы с собственным электромагнитным полем. Р. П. рассчитывают по методу теории возмущений, представляя их в виде ряда по степеням постоянной тонкой структуры



\[

\alpha=e^{2} / \hbar c \approx 1 / 137

\]



(где $e$ - заряд электрона), $n$-й член к-рого можно рассматривать как результат испускания и поглощения $\boldsymbol{n}$ виртуальных фотонов или электрон-позитронных пар. При вычислении Р. п. используется процедура перенормировки массы и заряда частицы.



Наибольший интерес представляют Р. П. к магнитным моментам электрона и мюона, к сверхтонкому расщеплению атомных уровней, радиационное смещение атомных уровней энергии, Р. П. к сечениям рассеяния электрона электроном или атомным ядром. Результаты расчетов Р. п. вплоть до величин 3-го порядка по степеням $\alpha$ блестяще согласуются с экспериментальными данными и свидетельствуют о справедливости квантовой электродинамики по крайней мере на расстояниях, бо́льших $10^{-15} \mathrm{~cm}$. Р. п. растут с ростом энергии, и эффективным параметром разложения (эффективным зарядом) при высоких энергиях являет- ся величина $\alpha \ln (\mathscr{E} / m)$, где $\mathscr{E}-$ энергия частицы в системе центра инерции, $m$ - ее масса в энергетич. единицах.



Р. п. могут быть в ряде случаев подсчитаны не только для электродинамич. процессов, но и для процессов, вызванных другими взаимодействиями. Напр., в квантовой хромодинамике вычисляются Р. п. к сечениям глубоко-неупругих процессов или к вероятностям распада мезонов со скрытым «очарованием».



При вычислении Р. П. к электродинамич. величинам с точностью выше 3-го порядка существенный вклад вносят процессы виртуального рождения адронов и эффекты слабого взаимодействия.



Лит.: [1] Ф е й н м а н Р., Теория фундаментальных процессов, пер. с англ., М., 1978; [2] Б̈ е р к е н Дж., Дрелл С. Д., Релятивистская квантовая теория, пер. с англ., т. 1-2, М., 1978. Б. Л. Ноффе.



РАДИУС СХОДИМОСТИ - см. Степенной ряд.



PAДИУС-ВЕКТОР - вектор, идущий в точку пространства из некоторой, заранее фиксированной точки.



PAЗБИЕНИЕ на под с истемы - способ разделения системы $N$ частиц на $K$ непересекающихся подсистем. Подробно Р. на подсистемы описывается указанием входящих в него подсистем. Напр., (132) (4) (65) обозначают разбиение системы шести частиц на три подсистемы. Порядок расположения подсистем и порядок частиц в обозначениях последних несуществен. Р. на подсистемы применяется в Квантовой теории рассеяния для описания и классификации состояний системы $N$ частиц.



Jum.: [1] We i nberg S., «Phys. Rev.», 1964, v. 133, № 1B, p. 232-56; [2] Якубовский О.А., «Ядерная физика», 1967, т. 5, с. 1312-20. С.П. Меркурьев, С. Л. Яковлев.



PAЗГPУกПИРОВКА ФОТОНОВ - см. Антигруппировка фотонов.



РАЗПЕЛЕНИЕ ДВИХЕНИЙ - ОДИН ИЗ ОСНОВНЫХ ПОДХОдов к исследованию систем дифференциальных уравнений с быстрыми и медленными переменными



\[

\dot{x}=f(x, y, \varepsilon), \dot{y}=\varepsilon g(x, y, \varepsilon), 0<\varepsilon \ll 1 .

\]



При этом подходе ищется близкая к тождественной замена переменных $x, y \rightarrow \xi, \eta$, к-рая позволила бы отделить уравнение для медленных переменных, то есть записать исходную систему в виде:



\[

\dot{\xi}=F(\xi, \eta, \varepsilon), \dot{\eta}=\varepsilon G(\eta, \varepsilon) .

\]



Часто дополнительно требуется, чтобы правая часть уравнения для быстрых переменных также зависела только от медленных переменных: $F=F(\eta, \varepsilon)$.



Замена переменных, осуществляющая Р. Д., ищется в виде ряда по малому параметру є. Точное Р. д. обычно невозможно, и ряды для замены переменных расходятся. Однако, обрывая эти ряды, можно осуществить Р. д. приближенно. Отбрасывая в уравнениях для новых переменных члены высокого порядка, зависящие от быстрых переменных, получают систему для приближенного описания движения, к-рая существенно проще, чем исходная. Примером Р. д. является использование процедур исключения быстрых фаз в усреднения методе.



Иногда приближенное Р. Д. удается осуществить не во всем фазовом пространстве, а лишь на нек-рой поверхности, как это делается в теории релаксационньх колебаний.



A. И. Нейumaдm.



РАЗДУВАЮШАЯСЯ ВСЕЛЕННАЯ - см. Квантовая теория гравитации, Космологические модели.



РАЗЛОХЕНИЕ ЕДИНИЧЫ - обобщение понятия спектральной меры, используемое в спектральном анализе несамосопряженных операторов и в теории квантового измерения. Пусть $\boldsymbol{H}$ - гильбертово пространство, $\boldsymbol{X}-$ измеримое пространство с $\sigma$-алгеброй измеримых подмножеств $\mathscr{B}(\mathfrak{X})$. Функция множеств $\mathscr{\mu}(B), B \in \mathscr{B}(\mathfrak{X})$, значениями к-рой являются ограниченные положительные операторы в $H$, называется разложением е диницы, если: 1) $\boldsymbol{\mu}(\mathfrak{X})=I$, где $I$ - единичный оператор в $H ; 2)$ для любого конечного или счетного набора $\left\{B_{i}\right\}$ попарно непересекающихся измеримых подмножеств



\[

\mathscr{\mu}\left(U_{i} B_{i}\right)=\Sigma_{i} \mu\left(B_{i}\right)

\]





\end{document}